\begin{frame}{Different Levels of Tokenization}
  \begin{itemize}
    \item[\textcolor{red}{$\blacktriangleright$}] \textbf{Character-level}
    \begin{itemize}
      \item Fine-grained, handles unknowns
      \item Long sequences, slow training
    \end{itemize}
    \vspace{0.7em}
    \item[\textcolor{red}{$\blacktriangleright$}] \textbf{Word-level}
    \begin{itemize}
      \item Intuitive, natural boundaries
      \item Large vocab, OOV (Out-of-Vocab) issues
    \end{itemize}
    \vspace{0.7em}
    \item[\textcolor{red}{$\blacktriangleright$}] \textbf{Subword-level}
    \begin{itemize}
      \item Balance of generalization \& compactness
      \item Requires segmentation algorithm
    \end{itemize}
    \vspace{0.7em}
    \item[\textcolor{red}{$\blacktriangleright$}] \textbf{SentencePiece/Unigram LM}
    \begin{itemize}
      \item Learned from data, language-agnostic
    \end{itemize}
  \end{itemize}
\end{frame}


\begin{frame}{Different Levels of Tokenization}
  \centering
  \includegraphics[width=0.60\paperwidth,height=0.60\paperheight,keepaspectratio]{assets/Lecture-7_Large_Language_Models/pdf_page_050.png}
\end{frame}

\begin{frame}{Different Levels of Tokenization}
  \centering
  \includegraphics[width=0.60\paperwidth,height=0.60\paperheight,keepaspectratio]{assets/Lecture-7_Large_Language_Models/pdf_page_051.png}
\end{frame}
